% ============================================
% Johns Hopkins Letter Template
% ============================================

\documentclass[11pt,letterpaper]{letter}

\usepackage{graphicx}
\usepackage{eso-pic}
\usepackage{geometry}
\usepackage{microtype}
\usepackage[utf8]{inputenc}
\usepackage[hidelinks]{hyperref}

\geometry{left=25mm,right=25mm,top=25mm,bottom=25mm}

\newcommand{\PutJHULogoFG}[1]{%
  \AddToShipoutPictureFG*{%
    \AtPageUpperLeft{%
      \hspace{10mm}%
      \raisebox{-38mm}{%
        \includegraphics[width=6cm]{#1}%
      }%
    }%
  }%
}

\signature{Decory Edwards}
\address{
Department of Economics \\
Johns Hopkins University \\
Baltimore, MD 21218 \\
\texttt{dedwar65@jh.edu}
}

\begin{document}
\PutJHULogoFG{Images/jhulogo.png}

\begin{letter}{
Search Committee \\
Department of Economics \\
Hamilton College \\
Clinton, NY
}

\opening{Dear Members of the Search Committee,}

I am writing to apply for the Visiting Assistant Professor position in Economics at Hamilton College. I am a Ph.D.\ candidate in the Department of Economics at Johns Hopkins University (advisors: Christopher D.\ Carroll, Nicholas W.\ Papageorge, and Stelios Fourakis), and I expect to receive my degree in May 2026. My primary specializations are macroeconomic modeling, wealth and income dynamics, and computational economics. I am committed to excellence in undergraduate teaching and to fostering an inclusive classroom environment that supports students from diverse backgrounds.

My research agenda focuses on the ability of macroeconomic models to generate empirically realistic distributions of wealth and income over the life cycle. A key driver of that work is modeling heterogeneity in the rate of return on assets, and understanding how return dispersion contributes to portfolio accumulation across households.

In my job market paper, I calibrate a life cycle heterogeneous agent model to match wealth moments from the Survey of Consumer Finances by allowing households to differ in their portfolio returns. The model helps explain observed wealth concentration and return dispersion. In a related research project using the Health and Retirement Study, I examine how trust influences both income growth and asset returns. I find a hump shaped relationship for trust and returns (consistent with the literature) and characterize the distribution of the persistent component of returns to net worth. These experiences have equipped me to frame research strategies and design modeling approaches, execute econometric and simulation analysis with large micro data sets, and translate results into clear and rigorous written deliverables for both academic and practitioner audiences.

Hamilton’s open curriculum and emphasis on intellectual curiosity are especially appealing to me. In such an environment, students encounter economics not as a checklist of required courses, but as a set of analytical tools for understanding inequality, markets, institutions, and public policy. My research on wealth concentration, return heterogeneity, and fiscal design provides rich material for engaging students in discussions about distribution, opportunity, and economic decision making. In a five course teaching load, I would welcome the opportunity to teach introductory and intermediate courses while offering electives that integrate macroeconomics, public finance, and quantitative reasoning. I strive to create classrooms that are rigorous yet accessible, where students feel encouraged to participate, challenge ideas respectfully, and develop confidence in their analytical skills.

I am enthusiastic about relocating to central New York and engaging fully in Hamilton’s residential liberal arts community. The College’s commitment to diversity, need blind admissions, and meeting full demonstrated financial need reflects a deep institutional investment in access and equity. I would seek to further these goals by mentoring students from a broad range of academic and socioeconomic backgrounds and by designing courses that connect economic theory to the lived experiences of diverse communities.

I believe I would be a strong fit for this role because:
\begin{enumerate}
    \item My research on inequality and wealth dynamics complements core themes in undergraduate economics education.
    \item I am committed to inclusive, discussion driven, and analytically rigorous teaching.
    \item I value close faculty student engagement and the collaborative intellectual culture that defines Hamilton College.
\end{enumerate}

Thank you very much for your consideration. I would welcome the opportunity to contribute to the Department of Economics at Hamilton College.

\closing{Sincerely,}

\end{letter}
\end{document}
